% 相互奇异：(a) [0,1]上的长度与点1/2处质量1/2的点测度，N={1/2}；(b) 两组互不相交的稠密可数点：有理点1/2,1/3,2/3,1/4,3/4（蓝）与sqrt2-1≈0.414, sqrt2-5/4≈0.164, sqrt2-7/10≈0.714, sqrt2-6/5≈0.214, sqrt2-1/2≈0.914（绿），高度仅示意权重递减。横轴1单位=4.2cm。
\begin{tikzpicture}[
  lab/.style={font=\figurefont\fontsize{8}{10}\selectfont,text=FigureInk},
  note/.style={font=\FigureNoteFont\fontsize{7.5}{9}\selectfont,text=FigureMuted,align=center},
  ax/.style={draw=FigureStroke,line width=.75pt,-{Stealth[length=2.2mm]}},
  spA/.style={draw=FigureBlue,line width=1.1pt},
  spB/.style={draw=FigureGreen,line width=1.1pt,dashed}]
\def\u{4.2}
\begin{scope}
  \draw[draw=FigureBlue,fill=FigureBlueFill,line width=.7pt] (0,0) rectangle (\u,.6);
  \draw[ax] (-.3,0)--(\u+.5,0);
  \draw[draw=FigureGreen,line width=1.4pt,-{Circle[length=1.6mm]}] ({\u/2},0)--({\u/2},1.4);
  \node[lab] at ({\u/4},.95) {$\lambda|_{[0,1]}$};
  \node[lab,right] at ({\u/2+.1},1.45) {$\tfrac12\delta_{1/2}$};
  \foreach \x/\t in {0/0,0.5/\frac12,1/1}{\draw[FigureStroke,line width=.6pt] ({\u*\x},0)--({\u*\x},-.08); \node[lab,below] at ({\u*\x},-.06) {$\t$};}
  \node[note] at ({\u/2},-.95) {$N=\{1/2\}$：长度为零，点质量全在其中};
  \node[lab] at ({\u/2},-1.5) {(a) 长度与点质量};
\end{scope}
\begin{scope}[xshift=5.6cm]
  \draw[ax] (-.3,0)--(\u+.5,0);
  \foreach \x/\h in {0.5/1.4,0.3333/1.05,0.6667/0.8,0.25/0.6,0.75/0.45}{\draw[spA] ({\u*\x},0)--({\u*\x},\h);}
  \foreach \x/\h in {0.4142/1.3,0.1642/1.0,0.7142/0.75,0.2142/0.55,0.9142/0.4}{\draw[spB] ({\u*\x},0)--({\u*\x},\h);}
  \foreach \x/\t in {0/0,1/1}{\draw[FigureStroke,line width=.6pt] ({\u*\x},0)--({\u*\x},-.08); \node[lab,below] at ({\u*\x},-.06) {$\t$};}
  \draw[spA] (.2,1.95)--(.55,1.95); \node[lab,right] at (.55,1.95) {有理点};
  \draw[spB] (2.2,1.95)--(2.55,1.95); \node[lab,right] at (2.55,1.95) {$\mathbb Q+\sqrt2$};
  \node[note] at ({\u/2},-.95) {两组点都稠密，但互不相交};
  \node[lab] at ({\u/2},-1.5) {(b) 交错的稠密点质量};
\end{scope}
\end{tikzpicture}
